\section{Introduction}


The blockchain industry stands at an impasse. Public blockchains
(Bitcoin, Ethereum) offer transparency and permissionless access but
cannot issue regulated financial instruments without violating
securities laws. Permissioned chains (Hyperledger, R3 Corda) satisfy
regulators but abandon the neutrality and composability that make
public blockchains valuable. Centralized exchanges (MtGox, Bitfinex,
Coincheck) combine the worst of both: regulatory exposure without
user custody~\cite{hacks}. Users have lost more than 980{,}000 BTC
(approximately USD 40 billion at current rates) to exchange hacks
since 2011.

Lux is designed to resolve this impasse. By operating under an
Isle of Man money transmitter license while retaining non-custodial
key management, open validator participation (subject to KYC/AML),
and public transaction auditability, Lux offers a blockchain that is
\emph{simultaneously} regulated and decentralized. The network's
hundred-validator founding set is both small enough to operate under
a coherent compliance regime and large enough to resist the
centralization risks that have plagued purely federated designs.

Lux addresses three additional concerns that existing blockchains do
not handle well:

\begin{itemize}[nosep]
  \item \textbf{Quantum resistance}: Shor's algorithm breaks ECDSA and
        RSA on a sufficiently large quantum computer. Lux adopts
        lattice-based primitives (from the \texttt{luxfi/lattice}
        library, based on CKKS~\cite{ckks} and BFV~\cite{bfv}) at
        every layer that touches user funds.
  \item \textbf{Regulatory compliance}: Lux integrates KYC/AML at the
        protocol layer via the Identity Chain (I-Chain) and the
        Decentralized Identifier (DID) specification. No unregistered
        security is issuable on the network.
  \item \textbf{Interoperability}: Lux bridges to existing blockchains
        through multi-party threshold signatures
        (MPC)~\cite{threshold-mpc}, providing cross-chain liquidity
        without trusted custodians.
\end{itemize}

%% ============================================================
